\documentclass[10pt,a4paper]{amsart}
\usepackage[margin=19mm]{geometry}
\usepackage{amsmath,amssymb,mathtools}
\usepackage[scr=rsfs,cal=euler]{mathalfa}
\usepackage{xfrac}
\usepackage[unicode,hidelinks,bookmarks=false]{hyperref}
\newcommand{\ie}{i.e.\ }
\newcommand{\cf}{cf.\ }
\newcommand{\resp}{respectively\ }
\newcommand{\Subsection}{\subsection}
\providecommand{\nobreakdash}{\nobreak\hbox{-}}
\setlength{\emergencystretch}{2em}
\multlinegap0pt
\usepackage{bm}
\usepackage{bbm}
\usepackage{dsfont}
\usepackage{xspace}
\usepackage{cancel}
\usepackage{mathtools}
\usepackage{amsthm}
\makeatletter
\@ifundefined{theorem}{\newtheorem{theorem}{Theorem}}{}
\@ifundefined{corollary}{\newtheorem{corollary}[theorem]{Corollary}}{}
\@ifundefined{lemma}{\newtheorem{lemma}[theorem]{Lemma}}{}
\makeatother
\DeclareMathOperator*{\sign}{sign\,}
\newcommand{\ar}{\mathbb{R}}
\newcommand{\beq}{\begin{equation}}
\newcommand{\eq}{\end{equation}}
\title[On the differentiability of the Sobolev norm of $W^{k,1}(M)$]{On the differentiability of the Sobolev norm of $W^{k,1}(M)$ on compact Riemannian manifolds}
\author{Duvan Cardona, Julio Delgado, Andres Felipe Mu\~noz-Tello}
\date{Source: arXiv:2606.15879v1 (CC BY 4.0)}
\begin{document}
\maketitle

\noindent\emph{Source and licence.} On the differentiability of the Sobolev norm of $W^{k,1}(M)$ on compact Riemannian manifolds. Version arXiv:2606.15879v1 (CC BY 4.0), distributed under the Creative Commons Attribution 4.0 International licence, as recorded in the arXiv licence metadata for this exact version and shown on the abstract page https://arxiv.org/abs/2606.15879v1. Full rights notice: licenses/arxiv-2606.15879-CC-BY-4.0.txt.\par\smallskip

\noindent\textit{Editorial scope.} Counted unit: the Theorem whose source label is \texttt{main1a} in the article (source0.tex lines 292--295, source label \texttt{main1a}), delivered together with the article's own complete original proof of it (source0.tex lines 297--332, one \texttt{proof} environment of 36 lines). No mathematics is written, repaired or re-derived: statement and proof are the article's own text line for line. Required context retained uncounted: ftja46 (lines 264--268); cordd2 (lines 272--277). Every definition, display and label of those lines is retained unchanged. Omitted from this package and not redistributed: 1--263, 269--271, 278--291, 296--296, 333--1051. Nothing inside a retained excerpt is omitted or altered: every retained body is delivered exactly as the article's lines are written, so no omission receipt of any kind accompanies this package. Citations: the retained excerpts carry no citation command at all, so no bibliography entry of the article is transcribed and no reference is redistributed. Reference adaptation: none was needed and none was made; every reference inside the delivered unit resolves inside it, so no unresolved reference or citation is produced. Printed slips kept literal: no printed word, symbol, hypothesis or number of the counted statement or of its proof was altered. Layout and class: the article's own class is \texttt{amsart} with packages including graphicx, amsmath, amssymb, amsthm, enumitem, bbm, nag, onlyamsmath, csquotes, url, latexsym, mathrsfs, delarray, amsfonts; the wrapper is the lane's pinned \texttt{amsart} editorial wrapper at 10pt on A4, which restates the theorem-like environments the retained text needs and adds the article's own macro definitions that the retained text needs (\texttt{\textbackslash ar}, \texttt{\textbackslash beq}, \texttt{\textbackslash eq}, \texttt{\textbackslash sign}), with extra wrapper packages (bm, bbm, dsfont, xspace, cancel, mathtools). The delivered numbering is the unit's own. Assets: no retained span contains a figure environment, a graphics-inclusion command, a PDF-page inclusion, a table or a TikZ picture; no image file is copied into this package, the package declares no asset dependency and no image file is redistributed with it. Licence: version arXiv:2606.15879v1 of this article is distributed under the Creative Commons Attribution 4.0 International licence, as recorded in the arXiv licence metadata for this exact version and shown on the abstract page https://arxiv.org/abs/2606.15879v1. A full rights notice is delivered at licenses/arxiv-2606.15879-CC-BY-4.0.txt. Characters: every retained excerpt is pure ASCII, so the delivered engine, XeLaTeX with its default Latin Modern text font, reports no missing character.
\begin{lemma}\label{lmain1a} Let $f,\, h:\Omega\rightarrow \ar$  be functions and $x\in\Omega$ such that $h(x)\neq 0$. Then 
\beq\lim\limits_{t\rightarrow 0}\frac{|f(x)+th(x)|-|f(x)|}{t},\label{j6nr}\eq
exists if and only if $f(x)\neq 0$. Moreover, if that is the case
\beq\lim\limits_{t\rightarrow 0}\frac{|f(x)+th(x)|-|f(x)|}{t}=\sign(f(x))h(x).\label{ftja46}\eq
\end{lemma}

\begin{corollary} \label{cordd2} Let $f, h:\Omega\rightarrow \ar$ be measurable functions such that $h(x)\neq 0$ $\mu$-a.e. Then 
\beq\lim\limits_{t\rightarrow 0}\frac{|f(x)+th(x)|-|f(x)|}{t} \mbox{ exists } \mu-a.e\label{limj57}\eq
if and only if $f(x)\neq 0$ $\mu$-a.e.\\

Moreover, if that is the case the formula \eqref{ftja46} holds.
\end{corollary}

\begin{theorem}\label{main1a} Let $M$ be a closed manifold. Let $B$ be the subset of $W^{k,1}(M)$ above defined. Then, the set of $G$-differentiability of the norm $\varphi:=\Vert\cdot\Vert_{W^{k,1}}$ is $B$. Moreover, the G\^ateaux derivative of  $\varphi$ at  $\tilde{f}\in B$ and the direction $h\in W^{k,1}(M)$ is given by 
	\beq\partial_{G}\varphi(\tilde{f})(h)=\sum_{\|\alpha\|_{\mathbb{N}^{n}}\leq k}\left[ \int_{M  } \sign(D^{\alpha}f(x))D^{\alpha}h(x)d\mu(x)\right] ,\label{Gder67}\eq
where $f$ is any representative of the class  $\tilde{f}$ and $D^{\alpha}f$, $D^{\alpha}h$ respectively are the derivatives of $f$ and $h$. 
	\end{theorem}

\begin{proof} %We will present the full proof in the case $n=1$ for simplicity. A careful look at its details will convince the reader ****\\

We assume $\tilde{f}\in B$ and that $f$ is a representative of the class $\tilde{f}$. We start by observing that for any $h\in  W^{k,1}(M)$ the right hand side of the  formula \eqref{Gder67} is independent of the representative in the class $\tilde{f}$. \\
	
We are now going to prove the $G$-differentiability of $\varphi$ at $\tilde{f}$. Let $h\in  W^{k,1}(M)\setminus\{0\}$, an application of the  Lebesgue dominated convergence Theorem gives us
{\setlength\arraycolsep{0pt}
\begin{align*}
	\lim\limits_{t\rightarrow 0}\frac{\varphi(f+th)-\varphi(f)}{t}=& \lim\limits_{t\rightarrow 0} \frac{\sum\limits_{\|\alpha\|_{\mathbb{N}^{n}}\leq k}  \|D^{\alpha}f(x)+tD^{\alpha}h(x)\|_{L^{1}(M)}-\sum\limits_{\|\alpha\|_{\mathbb{N}^{n}}\leq k}  \|D^{\alpha}f(x)\|_{L^{1}(M)}}{t}\\
=&\sum\limits_{\|\alpha\|_{\mathbb{N}^{n}}\leq k}\lim\limits_{t\rightarrow 0}\frac{  \|D^{\alpha}f(x)+tD^{\alpha}h(x)\|_{L^{1}(M)}-  \|D^{\alpha}f(x)\|_{L^{1}(M)}}{t}\\
=&\ \sum\limits_{\|\alpha\|_{\mathbb{N}^{n}}\leq k}\left[  \int\limits_{M}\lim\limits_{t\rightarrow 0}\frac{  |D^{\alpha}f(x)+tD^{\alpha}h(x)|-|D^{\alpha}f(x)|}{t}d\mu(x)\right].
	 	\end{align*}} 
Since $D^{\alpha}h(x)\neq 0$ $\mu$-a.e. and  $D^{\alpha}f(x)\neq 0$ $\mu$-a.e., by applying Corollary \ref{cordd2} we have

	\begin{align*}
\lim\limits_{t\rightarrow 0}\frac{  |D^{\alpha}f(x)+tD^{\alpha}h(x)|-|D^{\alpha}f(x)|}{t}=\sign(D^{\alpha}f(x))D^{\alpha}h(x)
	\end{align*}
$\mu$-a.e., and

 \[\lim\limits_{t\rightarrow 0}\frac{\varphi(f+th)-\varphi(f)}{t}=\sum\limits_{\|\alpha\|_{\mathbb{N}^{n}}\leq k}\left[ \int_{M  } \sign(D^{\alpha}f(x))D^{\alpha}h(x) d\mu(x)\right].\]
	
We now assume that $\tilde{f}\in B^{c}$ and we will show that $\varphi$ does not have $G$-derivative at  $\tilde{f}$. For the set 
\begin{equation}\label{Zeq} Z=\{x\in M:f(x)=0\}\subset B^{c}, 
 \end{equation}
we have
$\mu(Z)>0.$ Since the manifold is compact, we also have $\mu(Z)<\infty$.  We consider $h=1_{Z}$, thus $h\in W^{k,1}(M)\setminus 0$. Since $h=0$ on $Z^c$ and for $\|\alpha\|\neq 0$, $D^{\alpha}h=0$ on $M$, for $t\in\ar$ we have  $|f+th|-|f|=|f|-|f|=0$ on  $Z^c$ and $|D^{\alpha}f+tD^{\alpha}h|-|D^{\alpha}f|=|D^{\alpha}f|-|D^{\alpha}f|=0$ on $M$. Hence, for the calculation of the $G$-derivative of $\varphi$ at $f$ in the direction $h$, is enough to consider the integration on $Z$.
 For $t\neq 0$ we have 
{\setlength\arraycolsep{0pt}
		\begin{align*}
		\frac{\varphi(f+th)-\varphi(f)}{t}=& \sum\limits_{\|\alpha\|_{\mathbb{N}^{n}}\leq k}\frac{  \|D^{\alpha}f(x)+tD^{\alpha}h(x)\|_{L^{1}(M)}-  \|D^{\alpha}f(x)\|_{L^{1}(M)}}{t} \\ 
		 =&\frac{  \|f(x)+th(x)\|_{L^{1}(M)}-  \|f(x)\|_{L^{1}(M)}}{t}\nonumber\\&\quad +\sum\limits_{0<\|\alpha\|_{\mathbb{N}^{n}}\leq k}\frac{  \|D^{\alpha}f(x)+tD^{\alpha}h(x)\|_{L^{1}(M)}-  \|D^{\alpha}f(x)\|_{L^{1}(M)}}{t}\nonumber\\ 
		 =&\frac{ \int_{Z}|t|dx}{t}=\frac{|t|}{t}\mu(Z)=\sign(t)\mu(Z).\nonumber\\
		\end{align*}}
Hence,  we have $(\partial_G^+\varphi)f(h)=\mu(Z)$ and $(\partial_G^+\varphi)f(h)=-\mu(Z)$. Therefore  the $G$-derivative of $\varphi$ at $f$ does not exists in $B^c$, which shows that the set of $G$-differentiability of the norm $\varphi$ is  $B$.\\
 
It is clear that the formula \eqref{Gder67} defines a bounded linear operator from  $ W^{1,1}(M)$ into $\ar$ for every $\tilde{f}$ in $B$, and this concludes the proof of the theorem.
\end{proof}

\end{document}
